\begin{table}[!htbp]\centering
\caption{Cambios asociados con la exposición regional previa}\label{tab:short_regional}
\begin{threeparttable}\small
\setlength{\tabcolsep}{4pt}
\begin{tabular}{lrrrr}
\toprule
Resultado & Coeficiente & EE & $p$ & $p$ previo \\
\midrule
Salario horario (log) & 0.017 & 0.012 & 0.154 & 0.243 \\
Formalidad (pp) & -1.827 & 0.346 & $<0.001$ & 0.842 \\
Autoempleo (pp) & 1.003 & 0.400 & 0.019 & 0.202 \\
\bottomrule
\end{tabular}
\begin{tablenotes}\footnotesize
\item Coeficiente de período posterior por índice de Kaitz estandarizado, ecuación~\eqref{eq:short_regional}. Una unidad equivale a una desviación estándar entre los 25 departamentos. Misma ventana, muestras de resultados, controles individuales y ponderadores que en el Cuadro~\ref{tab:short_main}. El modelo regional no incluye el indicador binario de educación baja. Errores agrupados por departamento. Los valores $p$ previos corresponden a adelantos de exposición regional.
\end{tablenotes}
\end{threeparttable}\end{table}
